\documentclass[11pt]{article}

\usepackage{amsmath,amssymb,amsthm,mathtools}
\usepackage[margin=1in]{geometry}
\usepackage{microtype}
\usepackage[hidelinks]{hyperref}

\newtheorem{theorem}{Theorem}[section]
\newtheorem{lemma}[theorem]{Lemma}
\newtheorem{corollary}[theorem]{Corollary}
\newtheorem{remark}[theorem]{Remark}
\newtheorem{conjecture}[theorem]{Conjecture}

\newcommand{\dd}{\,\mathrm d}

\title{A Triangle with a Three-Edge Tail:\\
An Odd-Walk Domination Proof}
\author{}
\date{}

\begin{document}

\maketitle

\begin{abstract}
Let $R_3$ be a triangle with a path of length three attached at one triangle
vertex.  We prove that every graphon $W$ of edge density $p\geq1/2$ satisfies
\[
 t(R_3,W)\geq p^4(2p-1)=\Phi_{R_3}(p).
\]
The new input is the special path-density inequality
\[
 t(P_5,W)^3\geq t(P_3,W)^5,
\]
where $P_j$ denotes the path with $j$ edges.  This is the $(3,5)$ case of
the Erd\H{o}s--Simonovits odd-walk theorem proved by Blekherman and Raymond.
We give a complete transfer from finite hosts to graphons on arbitrary
probability spaces, then combine the path ratio with the pointwise rooted
triangle bound.  The result classifies the formerly open Atlas graph $102$.
\end{abstract}


%%%%%%%%%%%%%%%%%%%%%%%%%%%%%%%%%%%%%%%%%%%%%%%%%%%%%%%%%%%%%%%%%%%%%%%%%%%%%%%
\section{Graphon notation and the graph}
%%%%%%%%%%%%%%%%%%%%%%%%%%%%%%%%%%%%%%%%%%%%%%%%%%%%%%%%%%%%%%%%%%%%%%%%%%%%%%%

Let $(\Omega,\mu)$ be an arbitrary probability space and let
$W\colon\Omega^2\to[0,1]$ be a symmetric measurable graphon.  Write
\[
 p=t(K_2,W),\qquad
 d=T_W\mathbf1,\qquad
 A=T_Wd=T_W^2\mathbf1,\qquad
 B=T_WA=T_W^3\mathbf1.
\]
Define the rooted triangle density
\[
 \tau(x)=\int_{\Omega^2}
 W(x,y)W(x,z)W(y,z)\dd\mu(y)\dd\mu(z).
\]

Let $R_3$ have a triangle rooted at $o$ and three additional vertices
$u,v,w$, with the only additional edges $ou,uv,vw$.  Integrating the tail
from its free end toward the root gives
\begin{equation}
 t(R_3,W)=\int_\Omega\tau(x)B(x)\dd\mu(x).
 \label{eq:R3-density}
\end{equation}
It has six vertices and six edges.  Every vertex added along the tail has
$x-1$ colour choices after its parent is coloured, so
\[
 \chi_{R_3}(x)=x(x-1)^4(x-2),
 \qquad \chi(R_3)=3.
\]
Consequently
\begin{equation}
 \Phi_{R_3}(p)=p^4(2p-1).
 \label{eq:R3-target}
\end{equation}


%%%%%%%%%%%%%%%%%%%%%%%%%%%%%%%%%%%%%%%%%%%%%%%%%%%%%%%%%%%%%%%%%%%%%%%%%%%%%%%
\section{The two path inequalities on arbitrary graphons}
\label{sec:paths}
%%%%%%%%%%%%%%%%%%%%%%%%%%%%%%%%%%%%%%%%%%%%%%%%%%%%%%%%%%%%%%%%%%%%%%%%%%%%%%%

Let $P_j$ denote the path with $j$ edges and $j+1$ vertices.  We need two
standard walk inequalities in a form that applies to the probability-space
generality used throughout the catalogue.

\begin{lemma}[Blakley--Roy for $P_3$]
\label{lem:BR}
Every graphon of edge density $p$ satisfies
\[
 t(P_3,W)\geq p^3.
\]
\end{lemma}

\begin{proof}
Since $d(x)=0$ implies $W(x,y)=0$ for almost every $y$, define quotients by
$d(x)$ to be zero on that null support.  Almost everywhere on the support of
$W$,
\[
 W=(Wdd')^{1/3}(W/d)^{1/3}(W/d')^{1/3},
\]
where $d=d(x)$ and $d'=d(y)$.  H\"older gives
\[
 p^3\leq
 \left(\int_{\Omega^2}W(x,y)d(x)d(y)\dd\mu^2\right)
 \left(\int_{\Omega^2}\frac{W(x,y)}{d(x)}\dd\mu^2\right)
 \left(\int_{\Omega^2}\frac{W(x,y)}{d(y)}\dd\mu^2\right).
\]
Each quotient integral is at most one.  The first factor is exactly
$t(P_3,W)$.
\end{proof}

\begin{lemma}[The $(3,5)$ odd-walk inequality]
\label{lem:odd-walk}
Every graphon $W$ on every probability space satisfies
\[
 t(P_5,W)^3\geq t(P_3,W)^5.
\]
\end{lemma}

\begin{proof}
For finite simple host graphs this is the $t=3,k=5$ instance of the
Erd\H{o}s--Simonovits odd-walk theorem; we use the proof of Blekherman and
Raymond \cite{BlekhermanRaymond}.  We now justify the graphon transfer
without assuming that the underlying probability space is atomless or
standard Borel.

Fix a finite graph $F$.  For each $N$, sample independent points
$X_1,\ldots,X_N$ from $\mu$.  Conditional on these points, and independently
for every unordered pair $i<j$, place the edge $ij$ in a simple graph $G_N$
with probability $W(X_i,X_j)$.  Auxiliary independent uniform random
variables on $[0,1]$ realize these Bernoulli choices.

Let $T_N(F)=t(F,G_N)$.  Among the $N^{v(F)}$ maps
$V(F)\to[N]$, at most $\binom{v(F)}2N^{v(F)-1}$ are non-injective.  For an
injective map, all host edges used by distinct edges of $F$ are distinct;
conditional independence and then integration of the sampled points show
that its expected indicator is exactly $t(F,W)$.  Hence
\begin{equation}
 \left|\mathbb E T_N(F)-t(F,W)\right|
 \leq\frac{\binom{v(F)}2}{N}.
 \label{eq:mean-convergence}
\end{equation}

For the variance, expand $T_N(F)^2$ as a sum over two vertex maps.  If the
two image sets are disjoint, their indicators depend on disjoint sampled
points and disjoint auxiliary edge variables, and are independent.  The
number of ordered pairs of maps with intersecting images is at most
$v(F)^2N^{2v(F)-1}$.  Every covariance has absolute value at most one, so
\begin{equation}
 \operatorname{Var}T_N(F)\leq\frac{v(F)^2}{N}.
 \label{eq:variance-convergence}
\end{equation}
Equations \eqref{eq:mean-convergence}--\eqref{eq:variance-convergence} imply
$T_N(F)\to t(F,W)$ in probability.

Apply this simultaneously to $F=P_3,P_5$.  Every realization of the finite
simple graph $G_N$ satisfies
\[
 T_N(P_5)^3\geq T_N(P_3)^5.
\]
If the claimed limiting inequality were false, continuity would give an
open neighbourhood of $(t(P_3,W),t(P_5,W))$ on which the reverse strict
inequality holds.  Simultaneous convergence in probability places
$(T_N(P_3),T_N(P_5))$ in that neighbourhood with probability tending to
one, contradicting the finite inequality.  The graphon statement follows.
\end{proof}

Combining the preceding lemmas gives the ratio in precisely the form needed
below.

\begin{corollary}[Odd-path weighted mean]
\label{cor:path-ratio}
If $p>0$, put $Z=t(P_3,W)$.  Then $Z>0$ and
\[
 \frac{t(P_5,W)}{Z}\geq p^2.
\]
\end{corollary}

\begin{proof}
Lemma~\ref{lem:BR} gives $Z\geq p^3>0$.  Lemma~\ref{lem:odd-walk} gives
\[
 \frac{t(P_5,W)}Z\geq Z^{2/3}\geq p^2.
\]
\end{proof}


%%%%%%%%%%%%%%%%%%%%%%%%%%%%%%%%%%%%%%%%%%%%%%%%%%%%%%%%%%%%%%%%%%%%%%%%%%%%%%%
\section{The positive theorem}
\label{sec:positive}
%%%%%%%%%%%%%%%%%%%%%%%%%%%%%%%%%%%%%%%%%%%%%%%%%%%%%%%%%%%%%%%%%%%%%%%%%%%%%%%

\begin{lemma}[Pointwise rooted triangle]
\label{lem:rooted-triangle}
For almost every $x$,
\[
 \tau(x)\geq\max\{2A(x)-p,0\}.
\]
\end{lemma}

\begin{proof}
Apply $ab\geq a+b-1$ to $W(x,y),W(x,z)$, multiply by
$W(y,z)\geq0$, and integrate in $y,z$.  This gives
$\tau(x)\geq2A(x)-p$; also $\tau(x)\geq0$.
\end{proof}

\begin{theorem}[Triangle with a three-edge tail]
\label{thm:main}
For every graphon $W$ of edge density $p\in[1/2,1]$,
\[
 \boxed{
 t(R_3,W)\geq p^4(2p-1)=\Phi_{R_3}(p).
 }
\]
\end{theorem}

\begin{proof}
Put
\[
 Z=\int_\Omega B(x)\dd\mu(x).
\]
The operator definitions and self-adjointness of $T_W$ give
\begin{align}
 Z&=\langle\mathbf1,T_W^3\mathbf1\rangle=t(P_3,W),
 \label{eq:Z-path}\\
 \int_\Omega B(x)A(x)\dd\mu(x)
 &=\langle T_W^3\mathbf1,T_W^2\mathbf1\rangle
 =\langle\mathbf1,T_W^5\mathbf1\rangle=t(P_5,W).
 \label{eq:BA-path}
\end{align}
Since $p\geq1/2>0$, Lemma~\ref{lem:BR} gives $Z>0$.  Define the probability
measure
\[
 \dd\nu(x)=\frac{B(x)}Z\dd\mu(x).
\]
By \eqref{eq:Z-path}, \eqref{eq:BA-path}, and
Corollary~\ref{cor:path-ratio},
\begin{equation}
 \mathbb E_\nu A=\frac{t(P_5,W)}{t(P_3,W)}\geq p^2.
 \label{eq:A-mean}
\end{equation}

The function $f(a)=\max\{2a-p,0\}$ is convex and nondecreasing.  Using
\eqref{eq:R3-density}, Lemma~\ref{lem:rooted-triangle}, Jensen,
\eqref{eq:A-mean}, and $p^2\geq p/2$, we obtain
\begin{align*}
 t(R_3,W)
 &\geq Z\,\mathbb E_\nu f(A)
 \geq Z f(\mathbb E_\nu A)
 \geq Zf(p^2)
 \\
 &=Z(2p^2-p)
 \geq p^3\,p(2p-1)
 =p^4(2p-1).
\end{align*}
The last expression is the target in \eqref{eq:R3-target}.
\end{proof}

At $p=1/2$ the target is zero, and the proof remains valid without division
by $2p-1$.  At $p=1$, $W=1$ almost everywhere and both sides are one.


%%%%%%%%%%%%%%%%%%%%%%%%%%%%%%%%%%%%%%%%%%%%%%%%%%%%%%%%%%%%%%%%%%%%%%%%%%%%%%%
\section{Atlas consequence and formalization boundary}
%%%%%%%%%%%%%%%%%%%%%%%%%%%%%%%%%%%%%%%%%%%%%%%%%%%%%%%%%%%%%%%%%%%%%%%%%%%%%%%

\begin{corollary}[Atlas 102]
NetworkX Atlas graph $102$, with graph6 code \texttt{EJe?}, is positive.
It has order and size six, chromatic number three, and
\[
 \chi_{F_{102}}(x)=x(x-1)^4(x-2).
\]
For every graphon of edge density $p\in[1/2,1]$,
\[
 t(F_{102},W)\geq p^4(2p-1)=\Phi_{F_{102}}(p).
\]
\end{corollary}

\begin{proof}
Its NetworkX edge list is
\[
 04,\ 05,\ 12,\ 13,\ 23,\ 34.
\]
The vertices $1,2,3$ form a triangle, and the remaining edges form the
three-edge path $3-4-0-5$.  Thus Atlas $102$ is exactly $R_3$.
\end{proof}

The proof works with ordinary, non-injective homomorphism densities.  The
sampling transfer explicitly controls collisions and variance and therefore
does not silently replace homomorphism counts by injective counts.  It also
does not require the graphon probability space to have atoms, to be atomless,
or to be standard Borel.

For formalization, the genuinely new infrastructure is the special
odd-walk comparison in Lemma~\ref{lem:odd-walk}.  One may formalize either
the finite Blekherman--Raymond $(3,5)$ theorem together with the sampling
transfer, or a direct analytic proof of that special case.  After that lemma,
the graphon proof uses only operator-coordinate peeling, the already
formalized $P_3$ bound, the existing pointwise rooted-triangle inequality,
and scalar Jensen for the positive part of an affine function.

This argument applies specifically because a single three-edge tail produces
the consecutive odd path ratio $t(P_5)/t(P_3)$.  It does not prove arbitrary
pendant-tree closure, multiple three-edge tails, a three-edge broom, or tails
attached at several triangle vertices.  In particular it does not classify
Atlas $97$ or $100$.


%%%%%%%%%%%%%%%%%%%%%%%%%%%%%%%%%%%%%%%%%%%%%%%%%%%%%%%%%%%%%%%%%%%%%%%%%%%%%%%
\appendix
\section{A shorter route, and a smaller gap}
\label{sec:shorter}
%%%%%%%%%%%%%%%%%%%%%%%%%%%%%%%%%%%%%%%%%%%%%%%%%%%%%%%%%%%%%%%%%%%%%%%%%%%%%%%

This appendix was written while preparing the Lean formalization.  It records
two reductions of the argument above and one obstruction, and it changes
nothing about what is proved.  Throughout write
\[
 a_k=t(P_k,W)=\langle\mathbf1,T_W^k\mathbf1\rangle,
 \qquad a_0=1,\quad a_1=p .
\]

\subsection{The convexity step is not needed}

In the proof of Theorem~\ref{thm:main} the function $f(a)=\max\{2a-p,0\}$ is
only ever integrated against $B\geq0$.  Since $f(a)\geq2a-p$ pointwise,
\eqref{eq:R3-density} and Lemma~\ref{lem:rooted-triangle} already give
\begin{equation}
 t(R_3,W)=\int_\Omega\tau B\dd\mu
 \ \geq\ \int_\Omega(2A-p)B\dd\mu
 \ =\ 2\int_\Omega AB\dd\mu-p\int_\Omega B\dd\mu
 \ =\ 2a_5-p\,a_3 ,
 \label{eq:jensen-free}
\end{equation}
by \eqref{eq:Z-path} and \eqref{eq:BA-path}.  No auxiliary probability measure
$\nu$, no Jensen, and no convexity or monotonicity of $f$ is used.

\subsection{The odd-walk theorem is stronger than the theorem needs}

By \eqref{eq:jensen-free}, Theorem~\ref{thm:main} reduces to the scalar
statement $2a_5-p\,a_3\geq p^4(2p-1)$, equivalently
\begin{equation}
 2\bigl(a_5-p^5\bigr)\ \geq\ p\bigl(a_3-p^3\bigr).
 \label{eq:defect}
\end{equation}

\begin{lemma}
\label{lem:reduction}
Suppose $p\geq1/2$ and
\begin{equation}
 a_5\ \geq\ p^2a_3 .
 \label{eq:weak-ratio}
\end{equation}
Then $t(R_3,W)\geq p^4(2p-1)$.
\end{lemma}

\begin{proof}
By \eqref{eq:jensen-free} and \eqref{eq:weak-ratio},
$t(R_3,W)\geq2a_5-p\,a_3\geq(2p^2-p)a_3=p(2p-1)a_3$.
Since $2p-1\geq0$ and $a_3\geq p^3$ by Lemma~\ref{lem:BR}, this is at least
$p(2p-1)p^3=p^4(2p-1)$.
\end{proof}

Inequality \eqref{eq:weak-ratio} is strictly weaker than
Lemma~\ref{lem:odd-walk}: from $a_5^3\geq a_3^5$ and $a_3\geq p^3$,
\[
 a_5\ \geq\ a_3^{5/3}=a_3\cdot a_3^{2/3}\ \geq\ a_3p^2 .
\]
So the whole formalization gap is \eqref{eq:weak-ratio}, one inequality among
three walk densities, rather than the Erd\H{o}s--Simonovits odd-walk theorem
together with a sampling transfer.  Every other ingredient is already
mechanically checked: Lemma~\ref{lem:BR} is \texttt{pow\_three\_le\_pathIntegral},
Lemma~\ref{lem:rooted-triangle} is \texttt{rootedTriangle\_ge}, and
\eqref{eq:R3-density} follows from the assignment-integral peeling ladder.

\subsection{What \eqref{eq:weak-ratio} is not}

Two remarks delimit the search, both negative.

First, \eqref{eq:weak-ratio} is not a statement about the spectral measure
alone.  Let $\sigma$ be the spectral measure of $\mathbf1$ for $T_W$, so that
$a_k=\int\lambda^k\dd\sigma$ with $\sigma$ a probability measure.  The
inequality $\int\lambda^5\geq(\int\lambda)^2\int\lambda^3$ is false for general
probability measures: for
$\sigma=\tfrac1{17}\delta_{-1}+\tfrac{16}{17}\delta_{1/2}$,
\[
 \int\lambda^3\dd\sigma=\tfrac1{17}>0,
 \qquad
 \int\lambda^5\dd\sigma=-\tfrac1{34}<0 .
\]
Adding an atom of small mass at $+1$ makes $\sigma$ satisfy the
Perron--Frobenius constraint
$\operatorname{supp}\sigma\subseteq[-\lambda_{\max},\lambda_{\max}]$ with
$\sigma(\{\lambda_{\max}\})>0$, and still leaves
$\int\lambda^5<0<(\int\lambda)^2\int\lambda^3$.  Hence no argument using only
the moment data of $\sigma$ can succeed: nonnegativity of $W$ has to enter
somewhere other than through the spectrum.  What it gives for free is only
$a_{2k+1}=\langle T_W^k\mathbf1,T_W\,T_W^k\mathbf1\rangle\geq0$.

Second, four natural routes fail, each of them on the complete bipartite
graphon $W_s$ with parts of measure $s$ and $1-s$, for which
\[
 p=2s(1-s),\quad a_2=s(1-s),\quad a_3=2s^2(1-s)^2,\quad
 a_4=s^2(1-s)^2,\quad a_5=2s^3(1-s)^3 .
\]

\begin{enumerate}
\item \emph{H\"older in the Blakley--Roy pattern.}  Factoring
 $Wdd'=(WAA')^{1/3}\bigl(Wd^3d'^3/(AA')\bigr)^{1/3}W^{1/3}$ and applying
 H\"older with exponents $(3,3,3)$ gives $a_3^3\leq a_5\cdot J\cdot p$, where
 $J=\iint W\,(d^3/A)(d'^3/A')$.  One would need $J\leq p^3$.  On $W_s$,
 $J=2s^2(1-s)^2$ while $p^3=8s^3(1-s)^3$, so $J\leq p^3$ forces
 $(1-2s)^2\leq0$.
\item \emph{H\"older with exponents $(5/3,5/2)$.}  This reduces
 \eqref{eq:weak-ratio} to $\langle h,T_Wh\rangle\leq1$ for
 $h=d^{5/2}/A^{3/2}$.  On every regular graphon $\langle h,T_Wh\rangle=1$
 exactly, so the reduced statement is tight and no crude estimate closes it.
\item \emph{A positive-definite correction.}  With $u=A-pd=T_W(d-p)$,
 $\langle u,T_Wu\rangle=a_5-2pa_4+p^2a_3$.  Were this nonnegative,
 \eqref{eq:weak-ratio} would follow from $a_4\geq p\,a_3$, itself a
 consequence of $a_3\leq(a_2a_4)^{1/2}$, $a_4\geq a_2^2$ and $a_2\geq p^2$.
 But on $W_s$ one has $u=\mp s(1-s)(1-2s)$ on the two parts and
 $\langle u,T_Wu\rangle=-2s^3(1-s)^3(1-2s)^2<0$ for $s\neq1/2$.
\item \emph{Tilting.}  Writing $a_5=a_2^2p_A$ and $a_3=p^2p_d$, where $p_A$
 and $p_d$ are the edge densities of $W$ with respect to $A\dd\mu/a_2$ and
 $d\dd\mu/p$, inequality \eqref{eq:weak-ratio} would follow from $a_2\geq p^2$
 together with $p_A\geq p_d$.  The latter is $a_5p^2\geq a_3a_2^2$, which on
 $W_s$ reads $4s(1-s)\geq1$.
\end{enumerate}

\subsection{A stronger statement consistent with all numerical evidence}
\label{sec:supermult}

The following would give \eqref{eq:weak-ratio} through
$a_5\geq a_2a_3\geq p^2a_3$, and would also contain Lemma~\ref{lem:BR} through
$a_3\geq a_1a_2\geq a_1^3$.

\begin{conjecture}
\label{conj:supermult}
For every graphon $W$ and all $m,n\geq0$, \ $a_{m+n}\geq a_ma_n$.
\end{conjecture}

Equality holds throughout whenever $W$ is regular, since then $a_k=p^k$.  The
statement was tested on roughly $5\times10^6$ pseudorandom graphons --- uniform
$[0,1]$ kernels, Bernoulli $0/1$ kernels, stochastic block kernels and
hub-like kernels, on ground sets of size $2$ to $14$ --- for every pair with
$m+n\leq7$.  The minimum of $a_{m+n}-a_ma_n$ over that sample was zero to
machine precision, attained only at regular kernels.  In spectral terms
Conjecture~\ref{conj:supermult} is the correlation inequality
$\int\lambda^{m+n}\geq\int\lambda^m\int\lambda^n$, which by the first remark
above is false for general probability measures; a proof must therefore use
more about $W$ than its spectrum.  We have no proof.

\subsection{Consequence for the formalization}
\label{sec:consequence}

The Lean development splits cleanly at \eqref{eq:weak-ratio}.  Everything else
is available, so the statement
\[
 \text{$p\geq\tfrac12$ and $p^2\,t(P_3,W)\leq t(P_5,W)$}
 \ \Longrightarrow\
 p^4(2p-1)\leq t(R_3,W)
\]
is unconditional Lean mathematics, and Atlas $102$ becomes \texttt{verified}
exactly when \eqref{eq:weak-ratio} is discharged --- either by a direct
analytic proof, or by formalizing the finite Blekherman--Raymond $(3,5)$
theorem together with the sampling transfer of Lemma~\ref{lem:odd-walk}.  The
second route is much the heavier: it needs a product probability space, a
conditional Bernoulli construction, a second-moment estimate and convergence in
probability, none of which the catalogue currently has.

%%%%%%%%%%%%%%%%%%%%%%%%%%%%%%%%%%%%%%%%%%%%%%%%%%%%%%%%%%%%%%%%%%%%%%%%%%%%%%%
\section{What the Lean formalization actually proves}
%%%%%%%%%%%%%%%%%%%%%%%%%%%%%%%%%%%%%%%%%%%%%%%%%%%%%%%%%%%%%%%%%%%%%%%%%%%%%%%

Theorem~\ref{thm:main} is formalized unconditionally, as
\texttt{satisfiesLowerBound\_102} in
\texttt{lean/Taeyoung/Methods/OddWalk/Row102.lean}.  Atlas $102$ is
\texttt{verified}.  Two differences of substance.

\subsection*{Lemma~\ref{lem:odd-walk} is discarded, and its gap is closed}

Appendix~\ref{sec:shorter} leaves \eqref{eq:weak-ratio} open and weighs two
ways to obtain it, calling the second --- the finite Blekherman--Raymond
theorem plus the sampling transfer proved in Lemma~\ref{lem:odd-walk} ---
much the heavier, since it needs a product probability space, a conditional
Bernoulli construction, a second-moment estimate and convergence in
probability.

None of that is formalized.  The proof of Lemma~\ref{lem:odd-walk} given above
is not used at all: the $(3,5)$ inequality $a_3^5\leq a_5^3$ is instead proved
outright, by the entropy argument of \texttt{notes/blekherman\_raymond.tex}
\S2 run against relative entropy (see the appendix of that note for how).
\eqref{eq:weak-ratio} then follows in two lines, since $a_3\geq p^3$ gives
$(p^2a_3)^3=p^6a_3^3\leq a_3^5\leq a_5^3$.

\subsection*{The theorem is proved through Appendix~\ref{sec:shorter}, not
through Section~\ref{sec:positive}}

The auxiliary probability measure $\nu$, the Jensen step and the function
$f(a)=\max\{2a-p,0\}$ of Theorem~\ref{thm:main}'s proof do not occur.  What is
formalized is \eqref{eq:jensen-free} followed by Lemma~\ref{lem:reduction};
the full odd-walk inequality is used only through the weaker ratio
\eqref{eq:weak-ratio}.

\begin{thebibliography}{9}

\bibitem{BlekhermanRaymond}
G.~Blekherman and A.~Raymond,
\emph{A new proof of the Erd\H{o}s--Simonovits conjecture on walks},
Graphs and Combinatorics 39 (2023), Article 53;
\href{https://arxiv.org/abs/2009.10845}{arXiv:2009.10845}.

\end{thebibliography}

\end{document}
